\documentclass[HARVARD,Utopia2COL]{WileyNJDv5}
\usepackage{amsmath,amssymb,booktabs,graphicx}
\hypersetup{hidelinks}
\articletype{Research Article}
\journal{Journal of Futures Markets}
\raggedbottom

\begin{document}
\title{When Does Posterior Integration Change Option Prices? Evidence from WTI Average Price Options}
\author[1]{Heriberto Espino Montelongo}
\authormark{ESPINO MONTELONGO}
\titlemark{POSTERIOR INTEGRATION AND WTI OPTION PRICES}
\address[1]{\orgname{Universidad de las Am\'ericas Puebla}, \orgaddress{\state{Puebla}, \country{Mexico}}}
\corres{Corresponding author: Heriberto Espino Montelongo.}

\abstract[Abstract]{
Integrating uncertainty in estimated volatility can change an option's point value without measuring the full uncertainty around that value. We quantify this distinction for WTI Average Price Options and compare it with the effect of changing the volatility information supplied to the pricing model. Holding the posterior and pricing model fixed, the integration correction is locally proportional to posterior variance times volatility curvature, whereas posterior price dispersion is first order in posterior standard deviation. A synthetic experiment maps the correction across information levels, contract horizons, moneyness, and partial fixing. In the October-2026 WTI sample, the mean absolute integration correction is \$0.000745 per barrel, compared with a mean 95\% posterior model-price interval width of \$0.2733. On 2,381 later-date valuations, prior option-implied volatility inputs yield substantially lower errors than the evaluated historical estimators, including a horizon-matched GARCH benchmark. Same-contract lagged implied volatility reproduces the improvement without cross-sectional smoothing, and prior vanilla-WTI option information preserves it across instruments. The evidence identifies when integration changes a point price and shows that, in this sample, volatility information is the larger pricing margin even when posterior price uncertainty remains appreciable.
}
\keywords{Bayesian option pricing; parameter uncertainty; Asian options; WTI crude oil; implied volatility}
\maketitle
\section{Introduction}\label{sec:short_intro}

Option prices depend on parameters that are estimated rather than observed. A Bayesian analysis can account for this uncertainty by integrating the conditional option value over the posterior distribution instead of fixing all parameters at point estimates. This idea is established in Bayesian option pricing and forecasting \cite{guidolin2003,romboutsstentoft2014,guptareisinger2014}. The practical question is narrower: when does that integration materially change the option value?

Arithmetic-average Asian options provide a useful setting because their exposure to volatility changes with moneyness, maturity, and the fraction of the averaging period already fixed. The numerical literature has developed efficient methods for this path dependence, including geometric-average control variates and conditioning approximations \cite{kemnavorst1990,curran1994}. Commodity applications additionally require a term structure of futures prices and, in richer specifications, stochastic convenience yields, jumps, or stochastic volatility \cite{schwartz1997,shirayatakahashi2011,ewaldwu2023}. The objective here is not to introduce another Asian-option model. It is to isolate the effect of parameter uncertainty within a transparent pricing model and to compare that effect with the empirical importance of the volatility information supplied to the model.

The distinction matters because two different questions are easily conflated. The first is whether averaging the conditional option value over a posterior distribution changes the posterior mean price. The second is whether the volatility input itself is aligned with the volatility information embedded in option prices. These are conceptually separate. The first is a nonlinear-averaging problem. The second concerns model specification and information content.

Let \(C^Q(\sigma)\) be the risk-neutral option value conditional on volatility \(\sigma\), and let \(\bar\sigma=\mathbb E[\sigma\mid\mathcal D]\). The posterior-integrated and posterior-mean prices are
\begin{equation}
 C_{PI}=\mathbb E[C^Q(\sigma)\mid\mathcal D],
 \qquad
 C_{PM}=C^Q(\bar\sigma).
 \label{eq:short_pi_pm}
\end{equation}
A second-order expansion gives
\begin{equation}
 C_{PI}-C_{PM}
 \approx
 \frac12 C^{Q\prime\prime}(\bar\sigma)
 \operatorname{Var}(\sigma\mid\mathcal D).
 \label{eq:short_taylor}
\end{equation}
Posterior integration therefore matters when parameter uncertainty coincides with curvature of the option-pricing function. By contrast, the dispersion of the posterior distribution of model prices is locally first order in the posterior standard deviation. A small PI--PM difference need not imply little uncertainty in the theoretical option value.

The empirical application uses WTI Average Price Options (APOs), whose payoff depends on the arithmetic average of daily first-nearby WTI futures settlements. WTI is useful for this comparison because historical volatility can change sharply, the futures curve varies across valuation dates, and option-implied volatility is observable from both the APO market and a separate vanilla-WTI option market. The design uses only information available before each target date when evaluating out-of-sample pricing performance.

Three findings organize the paper. First, a large synthetic experiment confirms Equation~\eqref{eq:short_taylor}: posterior integration becomes important primarily in weak-information states with substantial volatility curvature and little of the averaging period already fixed. Second, the observed October-2026 WTI sample lies in a region where posterior-integrated and posterior-mean prices are almost identical, even though the posterior distribution of model prices remains materially dispersed. Third, prior option-implied volatility inputs produce substantially lower out-of-sample pricing errors than the historical-volatility benchmarks evaluated here. Rolling historical windows, EWMA, and a horizon-matched GARCH(1,1) forecast do not reproduce that performance. A same-contract lagged-IV benchmark shows that the result does not require a flexible cross-sectional smile, and a separate validation using vanilla-WTI options shows that it is not confined to volatility inferred from the APO family itself.

These results do not identify a pure structural difference between physical and risk-neutral volatility. Implied volatility under a parsimonious pricing model can absorb risk premia, stochastic volatility, jumps, omitted futures-curve dynamics, liquidity, and settlement conventions. The supported statement is therefore empirical and conditional: in this sample, prior option-implied volatility contains substantially more information for subsequent APO settlement pricing than the historical-volatility inputs considered here, while posterior integration of the historical-volatility distribution changes the point-price mean much less.

The remainder of the paper is organized around those three findings. Section~\ref{sec:short_method} develops the posterior-pricing mechanism. Section~\ref{sec:short_design} describes the WTI data and out-of-sample design. Section~\ref{sec:short_results} reports the synthetic and empirical results. Section~\ref{sec:short_robustness} summarizes robustness and interpretation, with implementation details and extended diagnostics reported in the online supplement. Section~\ref{sec:short_conclusion} concludes.

\section{Posterior pricing and the integration correction}\label{sec:short_method}

\subsection{Historical inference and conditional valuation}
Historical log returns follow the physical-measure benchmark
\begin{equation}
 R_i\sim\mathcal N\!\left[
 \left(\mu-\frac12\sigma^2\right)\Delta t,
 \sigma^2\Delta t\right],
 \label{eq:short_returns}
\end{equation}
where $\mu$ is annualized drift, $\sigma>0$ is annualized volatility, and $\Delta t$ is the observation interval in years. The baseline priors are
\begin{equation}
 \mu\sim\mathcal N(0,1),\qquad
 \sigma\sim\operatorname{InvGamma}(2,0.1).
 \label{eq:short_priors}
\end{equation}
The inverse-gamma prior is placed directly on volatility, with shape $2$ and scale $0.1$. Empirical sampling uses random-walk Metropolis in $(\mu,\log\sigma)$, including the transformation Jacobian. For the synthetic experiment, $\mu$ is marginalized analytically conditional on $\sigma$ and the remaining posterior is evaluated by deterministic quadrature.

Under the risk-neutral measure, futures maturity $u$ follows
\begin{equation}
 dF_t^u=\sigma F_t^u\,dW_t^Q.
 \label{eq:short_q}
\end{equation}
The same Brownian motion drives all maturities. This common-factor assumption fixes the cross-maturity dependence used to value the average. Historical drift does not enter the valuation. Using the historically inferred $\sigma$ in Equation~\eqref{eq:short_q} defines the historical-volatility benchmark; its empirical adequacy is evaluated separately from the effect of integrating its posterior.

For fixing dates $t_1,\ldots,t_M$, the settlement average is
\begin{equation}
 A_T=\frac1M\sum_{j=1}^M F^{(1)}_{t_j},
 \label{eq:short_average}
\end{equation}
where $F^{(1)}_{t_j}$ is the first-nearby CL futures settlement applicable on date $t_j$. At valuation date $t$, realized fixings enter deterministically and unresolved fixings use their applicable futures maturities initialized from the current CL curve. The conditional call value is
\begin{equation}
 C_t^Q(\sigma)=B(t,T)\,
 \mathbb E^Q[(A_T-K)^+\mid\mathcal F_t;\sigma],
 \label{eq:short_call_value}
\end{equation}
with strike $K$, discount factor $B(t,T)$, and valuation information $\mathcal F_t$. Put values replace $(A_T-K)^+$ with $(K-A_T)^+$. Below, $C^Q(\sigma)$ suppresses the fixed contract state and applies to either option type.

Historical PI and PM are evaluated by Monte Carlo. Curran's conditioning approximation \citep{curran1994} supplies deterministic prices for scalar historical and option-implied inputs. Numerical checks compare Monte Carlo, independently scrambled Sobol pricing, and Curran calculations. Pricing approximation and posterior integration are distinct effects; the reported pricing method is identified in each empirical table.

\subsection{Point-price correction versus price dispersion}
Writing $\delta=\sigma-\bar\sigma$, a Taylor expansion gives
\begin{align}
 C^Q(\sigma)={}&C^Q(\bar\sigma)+C^{Q\prime}(\bar\sigma)\delta\notag\\
 &+\tfrac12 C^{Q\prime\prime}(\bar\sigma)\delta^2+R_3(\sigma).
 \label{eq:short_expansion}
\end{align}
Because $\mathbb E[\delta\mid\mathcal D]=0$, averaging this expression yields Equation~\eqref{eq:short_taylor} plus $\mathbb E[R_3\mid\mathcal D]$. When the third derivative is bounded in absolute value by $M_3$ over the relevant range,
\begin{equation}
 |\mathbb E[R_3\mid\mathcal D]|
 \leq\frac{M_3}{6}\mathbb E[|\delta|^3\mid\mathcal D].
 \label{eq:short_remainder}
\end{equation}
The approximation is useful when the posterior is sufficiently concentrated and higher-order curvature is controlled. The integration correction can have either sign; positive volatility curvature is not assumed.

The same posterior induces conditional values $C^Q(\sigma)$ with dispersion
\begin{equation}
 \operatorname{SD}(C^Q(\sigma)\mid\mathcal D)
 \approx |C^{Q\prime}(\bar\sigma)|
 \sqrt{\operatorname{Var}(\sigma\mid\mathcal D)}.
 \label{eq:short_delta}
\end{equation}
With bounded derivatives, the integration correction is locally second order in posterior standard deviation, whereas price dispersion has a first-order term when volatility sensitivity is nonzero. Thus PI--PM measures a shift in the point value, not the full uncertainty around it. Posterior model-price intervals quantify uncertainty in this conditional valuation; they are not predictive intervals for future market settlements.

\subsection{Mapping information and contract state}
The synthetic experiment contains 175,000 independent histories: 5,000 for each combination of $n\in\{21,42,63,126,252,504,1260\}$ and $\sigma_0\in\{0.10,0.20,0.35,0.50,0.80\}$. The contract layer crosses four horizons from 21 to 252 days, seven moneyness values from 0.60 to 1.50, and six partial-fixing states. The forward level is normalized to \$100 per barrel. For each history and state, the experiment compares the PI--PM correction with posterior variance times local volatility curvature. The experiment measures the approximation's performance and the magnitude of the correction across these states; the expansion itself is a standard analytical device.

\section{Data and empirical design}\label{sec:short_design}

\subsection{WTI Average Price Options}

The empirical sample is built from individual Barchart histories for WTI Average Price Options. The raw archive contains 213 unique contracts and 11,179 option-date observations across expiries from September 2026 through June 2029. The market target is the Barchart end-of-day settlement field. Calls and puts are retained as available, and moneyness is recomputed by valuation date rather than assigned once at the contract level.

For each option-date observation, the pricing data reconstruct three objects: realized first-nearby CL settlements for fixing dates that have already occurred, the CL futures contracts applicable to the remaining fixing dates, and the contemporaneous futures term structure used to initialize those remaining fixings. This preserves the contract feature that the first-nearby futures contract can change during the averaging month.

Historical volatility is estimated from WTI futures returns available by the valuation date. The baseline uses the labelled continuous/front-month \texttt{CL=F} return series as the physical-measure proxy. A robustness check reconstructs a first-nearby return history contract by contract from CL settlements and excludes returns spanning contract switches.

\subsection{Out-of-sample volatility comparisons}

The main empirical experiment asks how different volatility inputs affect subsequent APO pricing when only earlier-date option information is available. Seven expiries pass the committed-data coverage requirement, producing 2,381 out-of-sample contract-date observations over 15 target dates.

The historical-volatility benchmark is the posterior-integrated price based on the expanding historical return sample. Four pre-specified recent historical alternatives use the most recent 63, 126, or 252 returns and an EWMA estimate with a 63-business-day half-life. A Gaussian GARCH(1,1) provides a stronger dynamic benchmark. For each target date, GARCH parameters are estimated only from earlier returns; conditional variances are then forecast across the remaining fixing horizon and collapsed to a constant annualized volatility that matches the implied variance of the unresolved arithmetic-average component under the maintained one-factor pricing model.

The option-implied specifications are also restricted to earlier dates. Contract-level APO implied volatilities are obtained by inverting the Curran pricing approximation. A previous-day smile uses the latest completed earlier APO date for the relevant expiry, while an expanding specification uses all earlier dates with recency weights. The fitted relation is a function of moneyness and option type. Same-day APO settlements never enter the target-date volatility estimate.

A separate persistence benchmark tests whether the improvement requires a fitted smile. For each contract with an earlier implied volatility observation, the latest strictly prior IV for that same contract is carried forward and the contract is repriced using the current futures curve, discount factor, realized fixings, and remaining fixing schedule. This comparison is available for 2,366 of the 2,381 observations.

Pricing performance is summarized by mean error, mean absolute error (MAE), and root mean squared error (RMSE). Dependence across contracts observed on the same valuation date is preserved with valuation-date cluster bootstrap resampling.

\subsection{Cross-instrument validation}

The APO experiment still obtains implied volatility from the same option family later used as the pricing target. An independent validation therefore uses standard monthly American WTI options on futures (CME product LO) obtained through Databento. Official LO option settlements are matched to CL futures settlements, and contract-level implied volatilities are recovered with a Cox--Ross--Rubinstein futures-option tree.

For each APO target date, only LO observations from earlier dates are used. Implied-volatility smiles are fitted separately for the underlying CL futures contracts needed by the October-2026 APO fixing schedule. The resulting maturity-specific volatilities are mapped into the maintained scalar-volatility APO pricing model. The main comparison is restricted to 253 APO contract-date observations, spanning 10 target dates, that lie within the observed prior LO moneyness support for both external specifications.

The cross-instrument experiment is not intended to estimate a structural variance risk premium. Its role is narrower: it asks whether the out-of-sample improvement associated with prior option-implied volatility remains when the volatility information comes from a separate WTI option market rather than from earlier APO settlements.

\section{Results}\label{sec:short_results}

\subsection{Where the integration correction grows}
Figure~\ref{fig:short_mechanism} maps the synthetic evidence. Corrections are largest where historical information is weak and the pricing function has substantial volatility curvature. High volatility, long horizons, extreme moneyness, and limited realized fixing produce the strongest effects in the evaluated states. The largest absolute correction is \$0.1384 per barrel on the normalized \$100 forward; the largest cell-level 99th percentile is \$0.1271. These are observed simulation magnitudes, not bounds on all possible contract states.

The second-order approximation closely tracks the cell-level mean correction, with correlation 0.9995 and through-origin slopes close to one in the most adverse cells. This agreement links the measured correction to posterior variance and volatility curvature rather than to an unexplained difference between estimators.

Partial fixing generally reduces the correction as less of the average remains stochastic, but the relationship is not uniformly monotone. In Figure~\ref{fig:short_mechanism}C, the 252-day case initially increases before declining. Fixing additional observations changes both volatility exposure and the curvature of the residual option value. The integration effect therefore depends on the remaining payoff state, as well as on uncertainty in the volatility estimate.

\begin{figure*}[!tb]
\centering
\includegraphics[width=\textwidth]{figures/publication/fig01_mechanism_map.pdf}
\caption{Posterior integration across synthetic WTI-like states. Panel A shows posterior dispersion against historical sample size. Panels B and C show absolute PI--PM corrections by moneyness and partial fixing. Panel D compares cell-level mean corrections with Equation~\eqref{eq:short_taylor}; the correlation is 0.9995. Cell count records the number of scenario cells in each hexagonal bin.}\label{fig:short_mechanism}
\end{figure*}

\subsection{A small point correction with wider posterior price intervals}
Table~\ref{tab:short_october} reports the October-2026 baseline. PI and PM have nearly identical settlement-pricing errors. More directly, their mean absolute difference is \$0.000745 per barrel, while the mean central 95\% posterior model-price interval width is \$0.2733 per barrel. The interval width is about 367 times the correction, illustrating their different scales rather than defining a universal uncertainty ratio.

\begin{table}[!t]
\centering
\caption{October-2026 pricing and posterior model-price uncertainty, on 310 contract-date observations over 12 dates. Values are U.S. dollars per barrel. PI and PM use the same Monte Carlo pricing model; interval widths come from the historical-volatility posterior.}\label{tab:short_october}
\begin{tabular}{lr}
\toprule
Quantity & Value\\
\midrule
PI RMSE & 0.5177\\
PM plug-in RMSE & 0.5183\\
Mean absolute PI--PM correction & 0.000745\\
Mean 95\% posterior interval width & 0.2733\\
\bottomrule
\end{tabular}
\end{table}

The small correction follows the local mechanism: averaging a relatively concentrated posterior through this pricing function shifts its mean little. The wider price intervals reflect uncertainty that remains around the conditional value. Both historical-volatility point prices understate settlements on average in the October panel. Integrating the same posterior does little to remove that shared pricing discrepancy.

\begin{figure*}[!tb]
\centering
\includegraphics[width=\textwidth]{figures/publication/fig03_forward_q_cluster_bootstrap.pdf}
\caption{Pricing-error differences by APO expiry. Negative values indicate lower error for earlier-date APO implied volatility than for historical PI. Bars show 95\% percentile intervals from valuation-date cluster resampling. Volatility estimates use only earlier option dates; repricing conditions on the current futures curve and realized fixings.}\label{fig:short_forward_bootstrap}
\end{figure*}

\begin{table*}[!tb]
\centering
\caption{Earlier-date volatility information and later-date APO valuation. Panel A uses 2,381 exact common holdouts over 15 dates. Panel B uses the 2,366 observations with earlier IV for the same contract. MAE and RMSE are U.S. dollars per barrel. Historical PI is the canonical Monte Carlo baseline; scalar historical, GARCH, and option-implied inputs use Curran repricing.}\label{tab:short_forward}
\begin{tabular}{lrr}
\toprule
\multicolumn{3}{c}{Panel A: volatility inputs on the main holdouts}\\
Volatility input & MAE & RMSE\\
\midrule
Historical PI & 2.6187 & 3.4090\\
Rolling 63 returns & 4.5201 & 5.6465\\
Horizon-matched GARCH(1,1) & 2.6730 & 3.4441\\
Previous-day APO smile & 0.1075 & 0.1594\\
\midrule
\multicolumn{3}{c}{Panel B: same-contract persistence on exact common rows}\\
Volatility input & MAE & RMSE\\
\midrule
Historical PI & 2.6279 & 3.4173\\
Same-contract lagged IV & 0.0992 & 0.1403\\
Previous-day fitted APO smile & 0.1074 & 0.1594\\
\bottomrule
\end{tabular}
\end{table*}

\subsection{Volatility information dominates the evaluated historical comparators}
Table~\ref{tab:short_forward} compares volatility inputs on exact common samples. On the 2,381 holdouts, the historical PI baseline has RMSE 3.4090. The evaluated rolling estimators and EWMA do not improve on it; the strongest rolling comparator uses 63 returns. Horizon-matched GARCH has RMSE 3.4441, close to the expanding historical baseline. The full historical comparison appears in Section S1.

Prior APO option information produces substantially smaller errors. The previous-day smile has RMSE 0.1594, and the expanding earlier-date smile has RMSE 0.1725. Figure~\ref{fig:short_forward_bootstrap} shows improvements across all seven admitted expiries, with valuation-date bootstrap intervals below zero for the reported error differences.





Panel B explains part of that result. Carrying forward the same contract's earlier IV yields lower pooled error than fitting the previous-day smile. This descriptive ranking shows that the historical-versus-option-implied improvement does not require flexible cross-sectional smoothing. Persistence of option-implied states, combined with current-state repricing, is sufficient to reproduce it on these rows. The comparison does not establish general superiority of IV carry over fitted smiles.



\subsection{The option-implied improvement transfers across instruments}
Table~\ref{tab:short_lo} moves the volatility source outside the APO family. On 253 exact common-support observations, prior vanilla-WTI IV reduces RMSE from 0.5562 for historical PI to 0.2599. The corresponding prior APO smile has RMSE 0.2686. Date-cluster bootstrap intervals support the external improvement over historical volatility, while LO-versus-APO differences are smaller and are not uniformly distinguishable across MAE and RMSE.

\begin{table}[!t]
\centering
\caption{External volatility information on 253 October-2026 APO observations over 10 dates within prior LO moneyness support. Errors are U.S. dollars per barrel. Historical PI uses Monte Carlo; option-implied values use Curran. LO and APO volatility inputs use strictly earlier option dates.}\label{tab:short_lo}
\begin{tabular}{lrr}
\toprule
Volatility input & MAE & RMSE\\
\midrule
Historical PI & 0.4464 & 0.5562\\
Previous-day vanilla-WTI IV & 0.1559 & 0.2599\\
Previous-day APO IV & 0.1676 & 0.2686\\
\bottomrule
\end{tabular}
\end{table}

The external experiment follows the established commodity calibration logic of \citet{shirayatakahashi2011}, using earlier vanilla information in the present scalar model. Its role is to show that the improvement is not confined to reusing earlier prices from the target option family. Together with the persistence comparator, it identifies option-implied volatility as the larger empirical point-pricing margin in the evaluated samples.

\section{Robustness and scope}\label{sec:short_robustness}

Supporting checks address historical specification, numerical resolution, market activity, and external transfer. Online Supplement Section S1 reports all historical benchmarks, Section S2 covers numerical accuracy and posterior calibration, Section S3 assesses liquidity and date dependence, and Section S4 examines LO transfer and tree refinement.

Changing the physical-return model or return source leaves the central comparison intact. Standardized Student-$t$ innovations widen the volatility posterior and modestly improve historical-volatility pricing, while PI and PM remain close. Contract-by-contract reconstruction of the first-nearby CL return history changes the 2,381-row historical MAE/RMSE from 2.6187/3.4090 to 2.6468/3.4451. Thus the historical proxy's roll construction does not explain the observed advantage of option-implied inputs.

Numerical accuracy is important because the PI--PM corrections are small. High-precision Monte Carlo, independently scrambled Sobol pricing, and Curran conditioning recover representative corrections at a scale well above their stochastic numerical noise. They can differ in absolute prices by more than the correction itself, so agreement on PI--PM is checked separately from absolute price agreement. Simulation-based calibration tests the Gaussian posterior implementation. Repricing at posterior RMS volatility, $\sqrt{\mathbb E[\sigma^2\mid\mathcal D]}$, also changes pooled errors only slightly relative to posterior mean volatility under the same Curran approximation.

The LO result is insensitive to the evaluated scalar-transfer alternative. Matching the unresolved average's variance changes previous-day LO RMSE from 0.2599 to 0.2612 on the same 253 rows. An 80/160/320-step CRR audit shows practical stability of the production inversion. These checks support the reported comparison at its numerical scale; they do not establish an exact tree limit or an arbitrage-free fitted surface.

Market activity restricts interpretation. Many targets are end-of-day settlements with zero reported trading volume. The option-implied improvement persists in the pooled positive-volume sample and under open-interest restrictions, but transaction-active observations are sparse within some long-dated expiries. Date-cluster resampling preserves contemporaneous dependence among contracts; it does not eliminate serial dependence across neighboring dates. Fifteen target dates in the main forward exercise and ten in the LO comparison provide limited temporal evidence despite the larger number of contract-date rows.

The pricing model fixes constant volatility and one common futures factor. Implied volatility under this model can absorb risk premia, stochastic volatility, jumps, omitted curve dependence, liquidity, and settlement marking. The comparison therefore concerns the information supplied to a conditional valuation model, rather than an identified structural difference between physical and risk-neutral diffusion coefficients. Richer models could change both absolute fit and the sensitivity of prices to additional uncertain parameters. The present results identify the integration mechanism and its scale within the stated model and samples.

\section{Conclusion}\label{sec:short_conclusion}

This paper quantifies three consequences of volatility estimation for WTI Average Price Options: the point-price correction from posterior integration, the posterior dispersion of conditional prices, and the pricing effect of changing volatility information. Holding the posterior and pricing model fixed isolates the integration correction. Its local dependence on posterior variance and volatility curvature explains why it grows in weak-information, high-curvature states and why partial fixing can alter its magnitude.

In the observed October-2026 panel, integration changes the point price little even though posterior model-price intervals remain appreciably wider than that correction. Across later-date valuations, earlier option-implied volatility produces much lower settlement-pricing errors than the historical inputs evaluated here. Recent historical estimates and horizon-matched GARCH do not reproduce that performance; same-contract IV persistence and external vanilla-WTI information do. The empirical finding concerns volatility information under a fixed valuation model, rather than a structural volatility-risk-premium decomposition.

The resulting diagnostic extends the interpretation of small integration effects: closeness of PI and PM establishes little movement in the point value, not little uncertainty around it. For the evaluated WTI samples, the larger point-pricing effect comes from the volatility information supplied to the model. This conclusion remains conditional on the common-factor specification, available settlements, and limited number of market dates.


\bmsection*{Data availability}
The WTI APO and CL histories were obtained through a commercial Barchart subscription. The vanilla-WTI validation uses CME/NYMEX LO option and CL futures statistics accessed through Databento. Raw proprietary market-data files are not redistributed. The public repository provides code, query metadata, schemas, provenance records, diagnostics, hashes, and derived results consistent with the applicable data licenses.

\bmsection*{Code availability}
The Python package \texttt{bayesian\_asian\_options} contains tests, deterministic experiment entry points, and committed derived results. The public repository is \url{https://github.com/heri-espino/Bayesian-Uncertainty-in-WTI-APOs}. The reported computations correspond to Git commit \texttt{8d95e5313e7b764c0cb0e7f9ef06e54bcf5d75d7}.

\bmsection*{Conflict of interest}
The author declares no conflict of interest.
\bibliography{references}
\end{document}
